\documentclass[12pt]{article}
\usepackage[utf8]{inputenc}
\usepackage[T1]{fontenc}
\usepackage{hyperref}
\title{Buglere Language Resources for K-12 Curriculum Development}
\author{Compiled by Claude Sonnet 5.5 (claude-sonnet-5-5) for Andr\'as Kornai}
\date{}
\begin{document}
\maketitle

\section{Introduction}
Buglere (ISO 639-3 sab) is a Chibchan language of Panama, closely related to Guaym\'i, with about 18,000 speakers and two dialects, Sabanero and Bok\'ota. The search found only general overviews and a community-built online dictionary. It found no school materials, primers or dedicated grammars. This dossier is therefore short.

\section{Primary Educational Resources}
No primers, textbooks or school programmes for Buglere turned up in the search.

\section{Linguistic and Reference Materials}
\begin{itemize}
\item Wikipedia overview of Buglere: \url{https://en.wikipedia.org/wiki/Buglere}
\item Omniglot page on Buglere. It lists further references (a SIL archive page, a grammatical-relations paper, an 1896 word list). I did not open these, so they are unverified: \url{https://omniglot.com/writing/buglere.htm}
\end{itemize}

\section{Religious and Community Resources}
None found in the search.

\section{Digital and Audio Resources}
\begin{itemize}
\item Glosbe community dictionary entries for Buglere. They are mostly short and need checking before use: \url{https://glosbe.com/lib/sab}
\end{itemize}

\section{Curriculum Development Context}
Nothing found indicates that Buglere is taught in schools. A curriculum would need new materials built from fieldwork and from the unverified leads above. Further searches should target the SIL archive and the academic literature on Buglere.
\end{document}